\section{Definition files}
\label{sec:definitions}

\subsection{The format}

The subject leaves the format open and suggests XML. XML was adopted: the JDK ships a parser and
a schema validator, so the choice costs no dependency and brings structural validation for free.
\Cref{lst:definition} is a complete definition; it carries the six pieces of information required
by section 2.5 of the subject, in the notations the subject itself uses.

\begin{lstlisting}[style=xml, caption={A definition file using every feature of the format.},
  label={lst:definition}, captionpos=b]
<automaton name="Everything at once">
  <description>A 3D automaton that uses every feature of the format.</description>
  <dimension>3</dimension>                        <!-- dimensionality of the grid -->
  <size>40, 40, 40</size>                         <!-- one extent per dimension -->
  <boundary>periodic</boundary>                   <!-- extension; default: fixed -->
  <slice>(20, :, :)</slice>                       <!-- displayed plane, for 3D and more -->
  <neighborhoods>                                 <!-- additional neighborhoods -->
    <neighborhood name="G3">{(0, 0, 0), (1, 1, 0), (-1, 0, 0)}</neighborhood>
  </neighborhoods>
  <rule>                                          <!-- the evolution rule -->
    SI(EQ(COMPTER(G0), 1),
       SUPEQ(COMPTER(G26*), 13),
       EQ(COMPTER(G3), 2))
  </rule>
  <initial seed="42">RANDOM(50)</initial>         <!-- or a list of coordinates -->
  <colors>                                        <!-- extension -->
    <state value="0" color="#0F172A"/>
    <state value="1" color="#38BDF8"/>
  </colors>
</automaton>
\end{lstlisting}

The elements may come in any order. \code{boundary}, \code{slice}, \code{neighborhoods},
\code{colors} and \code{description} are optional. Coordinates and sets of coordinates are
written as in the subject; on a one-dimensional grid the parentheses may be dropped, so the
initial state of the Sierpi\'nski file is simply \code{(128)}.

\subsection{Three stages of validation}

A file written by hand will contain mistakes, and the worst way to learn about them is an
exception in the middle of a simulation. \code{XmlDefinitionReader} therefore accepts a file only
after three successive checks, each catching a different class of errors.

\begin{enumerate}
    \item \textbf{Structure.} The document is validated against an XML Schema bundled with the
          application (\code{automaton.xsd}): unknown or duplicated elements, missing mandatory ones,
          malformed colors or neighborhood names. These errors come with a line number.
    \item \textbf{Consistency.} The values are parsed and cross-checked: as many sizes as
          dimensions, a slice that fits the grid and leaves exactly two free axes, offsets and initial
          cells that have the right number of coordinates and lie inside the grid.
    \item \textbf{The rule.} The rule is compiled against the neighborhoods of the file. This proves
          that it is syntactically correct and only refers to neighborhoods that exist in the grid's
          dimension.
\end{enumerate}

A file that passes is guaranteed to be instantiable. A file that fails yields a
\code{DefinitionException}, a checked exception because every caller has to do something about
it, whose message starts with the name of the file. Twenty-seven kinds of invalid files are
exercised by a parameterized test.

\paragraph{Security.} A definition file may come from someone else. The XML parser is configured
to refuse document type declarations and external entities, which closes the classic XML external
entity attack; a test feeds it such a document. The size of a grid and the nesting depth of a
rule are bounded, so a hostile file cannot exhaust the memory or the stack.

\subsection{A definition is data}

The reader does not build an automaton. It builds an \code{AutomatonDefinition}, an immutable
record holding the sizes, the offsets of the neighborhoods, the rule \emph{as text} and the
initial state. Turning it into live objects is the separate job of \code{AutomatonFactory}:

\begin{lstlisting}
public static CellularAutomaton create(AutomatonDefinition definition) {
    ExpressionRule rule = compileRule(definition);
    Grid grid = new Grid(definition.shape(), definition.boundary());
    definition.initialState().applyTo(grid);
    return new CellularAutomaton(grid, rule);
}
\end{lstlisting}

This intermediate form is what the rest of the application manipulates. The menu displays
definitions without allocating any grid; a simulation instantiates one; the tests and the rule
explorer build definitions in code with a Builder, since the record has nine components:

\begin{lstlisting}
AutomatonDefinition glider = AutomatonDefinition.builder()
        .name("Glider")
        .size(20, 20)
        .boundary(BoundaryCondition.PERIODIC)
        .rule("SI(COMPTER(G0), OU(EQ(COMPTER(G8*), 2), EQ(COMPTER(G8*), 3)), EQ(COMPTER(G8*), 3))")
        .initialState(new InitialState.ExplicitCells(cells))
        .build();
\end{lstlisting}

The two forms of initial state allowed by the subject are the two implementations of a sealed
interface, \code{InitialState}: \code{ExplicitCells} and \code{RandomFill}. A sealed interface
states in the code that there are exactly two cases and lets the compiler verify that both are
handled. An \code{XmlDefinitionWriter} does the inverse of the reader, and a test checks that
writing then reading a definition gives back an equal definition.

\subsection{The files of a directory}

The menu required by the subject lists ``all the input files of the current directory''.
\code{DefinitionLibrary} scans a directory for XML files whose root element is
\code{<automaton>}, which keeps unrelated XML files (the \code{pom.xml} of the project, for one)
out of the menu. Each file becomes an \code{Entry}, another sealed interface: \code{Loaded} with
its definition, or \code{Broken} with the reason of the rejection. Invalid files are thus listed
with an explanation instead of silently disappearing.
